\documentclass{situationist-preprint}

\title{The Structure of Justification}
\msid{EPI-05-a}
\msfamily{EPI-05}
\msdiscipline{Epistemology}
\mstype{Principal manuscript}
\mslean{lean/Philosophy/EPI/EPI\_05.lean}
\mskeywords{regress problem; foundationalism; coherentism; infinitism; perceptual experience; support tree; Agrippa's trilemma}

\begin{document}
\maketitle

\begin{abstract}
If a belief is justified by other beliefs, and they by others, the chain of reasons either ends, circles back, or goes on forever. Foundationalism says that it ends in beliefs justified without support; coherentism that justification belongs to systems of mutually supporting beliefs; infinitism that the chain never ends. We argue for none of these. Justification has the structure of a tree whose roots are perceptual experiences. Every justified belief is supported: perceptual beliefs by experiences, other beliefs by inference from justified beliefs. Experiences are not beliefs, are neither justified nor unjustified, and so require no further support; the regress stops not at a privileged class of beliefs but where support passes from belief to experience. We prove that, if justification is generated by undefeated experiential grounding and transmitted by good inference, a belief is justified exactly when it has a support tree of finite height whose leaves rest on experiences; that no circle and no infinite chain of reasons generates justification; and that every justified empirical belief derives from some experience. A model shows that coherentist closure counts isolated systems as justified. We apply the result to memory, testimony, induction and expertise, place coherence where it belongs, and reply to the objections of Sellars, Davidson, BonJour, Klein and Haack. The results are formalized in Lean.
\end{abstract}

\section{The problem}

Suppose I believe that the kettle has boiled, and I am asked why. I say that I heard it click off. Asked why I believe that, I say that I heard a click from the kitchen and that the kettle is the only thing in the kitchen that clicks. Asked why I believe each of those, I give further answers. The questioning can be continued, and the question is what happens when it is. There are three possibilities, which together make up the trilemma that Sextus Empiricus attributed to Agrippa \citep[I.164--177]{Sextus}: the chain of reasons ends in something not supported by any further reason, which the sceptic calls an arbitrary hypothesis; it returns to a belief already used, which is circular; or it goes on without end, which is an infinite regress. The sceptic holds that none of the three yields justification, and concludes that no belief is justified.

The three classical responses each embrace one horn. \emph{Foundationalism} holds that the chain ends in basic beliefs, justified without being supported by other beliefs, and that every other justified belief is supported, ultimately, by basic ones. \emph{Coherentism} holds that justification is a property of systems of beliefs, and that a belief is justified by its membership in a coherent system, so that the apparent circularity of mutual support is not vicious. \emph{Infinitism} holds that the chain of reasons available to a justified believer is infinite and non-repeating.

We argue that none of these is correct, and that the trilemma rests on an assumption that should be rejected: that what supports a belief must itself be something that could be justified or unjustified. Our thesis is that justification has a tree structure whose roots are perceptual experiences requiring no further support. Every justified belief is supported by something. Inferentially justified beliefs are supported by other justified beliefs; perceptual beliefs are supported by the experiences on which they are based, in virtue of the seemings those experiences involve (EPI-19). Experiences are not beliefs. They make no claim that could be backed or undermined by reasons in the way beliefs can, and the question of their justification does not arise. The chain of support therefore ends, but not in a belief that stands without support; it ends where support passes from belief to experience. We shall say that the experiences are the roots of the tree: they are where it meets the ground, and they are not part of the structure that they hold up.

This view is not foundationalism in the classical sense, since it recognizes no justified belief that is unsupported and no belief that justifies itself. Nor is it the modest foundationalism that takes perceptual beliefs to be basic, for the same reason: a perceptual belief, on our view, is supported, by an experience, and is a node of the tree, not its root. It is not coherentism, since support is well founded and circles generate nothing. And it is not infinitism, since every tree is of finite height. The contribution of the present manuscript is to make the structure precise, to show that its central claims follow from two principles about the sources of justification and a third saying that there are no others, and to show how coherence, which coherentists rightly took to matter, fits within it.

\section{A history of the debate}\label{s:history}

\subsection{The ancient regress}

\citet[I.3]{AristotlePostAn} considered the regress for demonstrative knowledge and rejected both the view that demonstration goes on without end and the view that it is circular. He concluded that the first principles of demonstration are known without demonstration, by a different capacity. The structure he described, in which every demonstrated truth rests on finitely many indemonstrable first principles, is the template of foundationalism. Sextus turned the same alternatives against all claims to knowledge: the dogmatist must stop somewhere, and wherever he stops he assumes without proof something that his opponent may with equal right deny \citep[I.164--177]{Sextus}.

\subsection{Classical foundationalism}

Descartes sought foundations that could not be doubted, and found them in the existence of the thinking self and in clear and distinct perceptions \citep{Descartes1641}. The empiricist version, developed most fully by \citet{CILewis1946}, placed the foundations in beliefs about the given: reports of how things appear, which Lewis took to be certain because they make no claim beyond the experience they report. All other empirical beliefs were to be justified by probable inference from these. The programme failed on both sides. The supposedly certain beliefs about appearances are not certain: one can misdescribe one's experience, and the vocabulary in which appearances are reported is borrowed from the vocabulary of objects. And the inference from appearance-beliefs to beliefs about the world cannot be shown to be a good one without assuming the reliability of perception, which is what the inference was supposed to establish.

\subsection{The protocol sentence debate}

In the Vienna Circle, the question took the form of a debate about protocol sentences, the reports of observation on which science rests. \citet{Neurath1932} denied that protocol sentences have any privileged status: they are revisable like any others, and the scientist is like a sailor who must rebuild his ship on the open sea, never able to put it into dry dock. \citet{Schlick1934} replied that knowledge has a foundation in what he called affirmations, momentary confirmations of the form \enquote{here now blue}, which are not themselves part of science but the points at which science makes contact with reality. Schlick's affirmations, which are not hypotheses and cannot be revised because they do not last long enough to be written down, are the closest historical precursor of the view we defend: the contact point is not a belief in the system but an occurrence that the system answers to.

\subsection{The attack on the given}

\citet{Sellars1956} argued that the idea of a given, something that is both non-inferentially present to the mind and capable of justifying belief, is a myth. If the given has propositional content, it is the kind of thing that can be true or false and that needs justification; if it lacks propositional content, it cannot stand in logical relations to beliefs and so cannot justify them. \citet{Davidson1986} drew the coherentist conclusion: nothing can count as a reason for holding a belief except another belief, and experience is a cause of belief, not a reason for it. \citet{BonJour1978,BonJour1985} argued that a belief could be basic only if the believer had a reason to think that beliefs of its kind are likely to be true, which would make it non-basic, and developed a coherence theory on which justification is a matter of the inferential relations within the believer's whole system of beliefs. The fatal difficulty for these views was the \emph{isolation objection}: a coherent system of beliefs may be entirely cut off from experience, like a consistent novel, and yet coherentism must count its members as justified. BonJour's response was an observation requirement, that a coherent system must contain beliefs attributing reliability to spontaneous perceptual beliefs, which in effect readmits experience as an input; he later abandoned coherentism for a foundationalism of the given \citep{BonJourSosa2003}.

\subsection{Rafts, pyramids and crosswords}

\citet{Sosa1980} cast the debate as a contest between two images: the pyramid of foundationalism, in which every block rests on blocks below and ultimately on the ground, and the raft of coherentism, in which the planks hold one another up. He argued that both images share a formal conception of justification as supervening on logical relations among beliefs, and that a virtue-theoretic account cuts across them. \citet{Haack1993} offered a third image, the crossword: an entry is justified both by its clue, which corresponds to experience, and by its fit with intersecting entries, which corresponds to coherence. Her \enquote{foundherentism} allows experience to provide input without making any belief basic, and allows mutual support among beliefs. We accept the first half and reject the second, for reasons given in Section~\ref{s:coherence}.

\subsection{Infinitism}

\citet{Klein1999,Klein2007} revived the third horn. He argued that foundationalism makes the termination point arbitrary and coherentism makes justification circular, and that a belief is justified when the believer has available an infinite and non-repeating series of reasons for it, of which in practice she need produce only finitely many. \citet{Peijnenburg2007} and \citet{AtkinsonPeijnenburg2017} showed that an infinite chain of probabilistic support can determine a definite probability for its first member, so that infinite regresses are not formally incoherent. The companion manuscript EPI-05-b argues against infinitism in detail; Corollary~\ref{c:regress} below gives the structural reason it fails.

\section{Framework}\label{s:framework}

We distinguish two kinds of state. \emph{Beliefs} are the candidates for justification: states that represent their contents as true, are held for reasons or on grounds, and can be justified or unjustified. \emph{Experiences} are perceptual states in which things look, sound or feel a certain way to the subject. They have content, but they are not held for reasons, and the question whether they are justified does not arise for them, any more than it arises whether a sensation of warmth is justified.

\begin{definition}[Experiential grounding]\label{d:ground}
A belief $b$ is \emph{grounded in} an experience $e$ if $b$ is held on the basis of $e$, and $e$ makes it seem to the subject that the content of $b$ is true.
\end{definition}

\begin{definition}[Good inference]\label{d:infer}
A belief $b$ is held by a \emph{good inference} from premise beliefs $p_1, \ldots, p_k$ if it is based on them by a rule of inference whose use confers justification on its conclusion when the premises are justified.
\end{definition}

\begin{definition}[Defeat]\label{d:defeat}
A belief is \emph{defeated} if the subject has a defeater for it: a state that supports the falsity of its content, or undermines the connection between its grounds and its truth.
\end{definition}

Definition~\ref{d:infer} leaves open which rules are good. Deductive rules are, and EPI-03 argues that induction is, by a rule-circular argument that we return to in Section~\ref{s:cases}. Definition~\ref{d:defeat} leaves open what defeats. Defeaters may come from anywhere in the subject's system of beliefs, and this is one place where the holism the coherentist emphasized has its due.

\begin{assumption}[Experiential sufficiency]\label{a:exp}
A belief grounded in an experience, and not defeated, is justified.
\end{assumption}

\begin{assumption}[Transmission]\label{a:trans}
A belief held by a good inference from justified premises, and not defeated, is justified.
\end{assumption}

\begin{assumption}[No other source]\label{a:nos}
Justification is had only through the sources of Assumptions~\ref{a:exp} and~\ref{a:trans}: any property of beliefs that is possessed by every undefeated experientially grounded belief, and by every undefeated belief held by good inference from premises that possess it, is possessed by every justified belief.
\end{assumption}

\begin{assumption}[Empirical premises]\label{a:prem}
No empirical belief is held by inference from no premises.
\end{assumption}

Assumption~\ref{a:exp} is the conclusion of EPI-19: if it seems to a subject that $p$, she has prima facie justification for believing that $p$, which, undefeated, is justification. Assumption~\ref{a:trans} says that good inference passes justification from premises to conclusion. Assumption~\ref{a:nos} is the substantive structural premise. It says that justification does not come from nowhere: a belief is not justified by being part of a system, by being believed, or by having an infinite series of reasons available, unless that status is reached by the two sources. Formally, it says that the justified beliefs form the smallest class closed under the two sources, and that is what licenses proofs by induction on the generation of justification. Assumption~\ref{a:prem} restricts attention to empirical beliefs, for which premise-free inference is not available; we discuss the a priori in Section~\ref{s:limits}.

\begin{definition}[Support tree]\label{d:tree}
A belief $b$ has a \emph{support tree of height at most $n$} if either $b$ is grounded in an experience and is not defeated, or $n = m+1$ and $b$ is held by a good inference from premises each of which has a support tree of height at most $m$, and $b$ is not defeated.
\end{definition}

A support tree, drawn out, has the belief $b$ at its top, the premises of each inference as the children of its conclusion, and, at each leaf, an undefeated belief grounded in an experience. Below each leaf is the experience in which it is grounded. These are the roots. The tree may share nodes, since one premise may support several conclusions, and in that case the structure is a directed acyclic graph; unfolded, it is a tree, and we keep the term.

\section{Results}\label{s:results}

\begin{lemma}[Monotonicity]\label{l:mono}
If $b$ has a support tree of height at most $n$, it has one of height at most $n+k$, for every $k$.
\end{lemma}

\begin{proof}
By induction on the tree. A leaf is a support tree of any height. If $b$ is held by a good inference from premises with trees of height at most $m$, then by the induction hypothesis they have trees of height at most $m+1$, and $b$ has one of height at most $m+2$. The general case follows by induction on $k$.
\end{proof}

\begin{lemma}[Common bound]\label{l:bound}
If each of finitely many premises has a support tree of some height, they all have support trees of a common height.
\end{lemma}

\begin{proof}
By induction on the number of premises, taking as common height the sum of the heights, by Lemma~\ref{l:mono}.
\end{proof}

\begin{theorem}[Tree theorem]\label{t:tree}
A belief is justified if and only if it has a support tree of finite height.
\end{theorem}

\begin{proof}
If $b$ has a support tree, it is justified, by induction on the tree: the leaves are justified by Assumption~\ref{a:exp}, and each internal node by Assumption~\ref{a:trans}. Conversely, consider the property of having a support tree of some finite height. Every undefeated experientially grounded belief has it, with height $0$. If an undefeated belief is held by a good inference from premises each of which has it, then by Lemma~\ref{l:bound} the premises have trees of a common height $N$, and the belief has a tree of height $N+1$. By Assumption~\ref{a:nos}, every justified belief has the property.
\end{proof}

Theorem~\ref{t:tree} is the precise content of the claim that justification has a tree structure. It does not say merely that justification may be pictured as a tree; it says that being justified and having a finite well-founded support structure resting on experiences are the same thing. Lemma~\ref{l:bound} is where finiteness enters. A good inference has finitely many premises, and a finite set of finite trees has a finite bound. An infinitist must deny this, either by allowing inferences with infinitely many premises or by denying Assumption~\ref{a:nos}; we argue in Section~\ref{s:objections} that both denials are unmotivated.

\begin{theorem}[No rootless justification]\label{t:rootless}
Let $S$ be a set of beliefs, none of which is grounded in an experience, such that every good inference to a member of $S$ has a premise in $S$. Then no member of $S$ is justified.
\end{theorem}

\begin{proof}
Consider the property of not belonging to $S$. Every experientially grounded belief has it, since no member of $S$ is grounded. If a belief is held by good inference from premises none of which is in $S$, it is not in $S$, since otherwise one of its premises would be. By Assumption~\ref{a:nos}, every justified belief has the property.
\end{proof}

\begin{corollary}[Circles]\label{c:circle}
If finitely many beliefs, none grounded in experience, support one another in such a way that every good inference to any of them has a premise among them, none of them is justified.
\end{corollary}

\begin{corollary}[Regresses]\label{c:regress}
If $b_0, b_1, b_2, \ldots$ is an infinite chain of beliefs, none grounded in experience, such that every good inference to any $b_n$ has some $b_m$ among its premises, no $b_n$ is justified.
\end{corollary}

Both corollaries are instances of Theorem~\ref{t:rootless}, with $S$ the circle and the chain respectively. Together they dispose of two horns of Agrippa's trilemma as sources of justification. Neither mutual support nor endless support generates justification where there is no experiential input. Note what the corollaries do not say. They do not say that a belief cannot be supported by a circle or a regress in addition to an experiential tree; a belief that has a support tree is justified whatever else supports it. They say that circles and regresses add no justification that was not already there.

\begin{theorem}[Experiential dependence]\label{t:dep}
Every justified empirical belief derives, through a finite chain of good inferences, from some experience: either it is grounded in an experience, or it is held by good inference from a premise that so derives.
\end{theorem}

\begin{proof}
Consider the property of deriving from some experience. Every experientially grounded belief has it. If a belief is held by good inference from premises that have it, then, since by Assumption~\ref{a:prem} there is at least one premise, the belief has it too. By Assumption~\ref{a:nos}, every justified belief has it.
\end{proof}

Theorem~\ref{t:dep} is the empiricist core of the view: no empirical belief is justified independently of experience. It depends on Assumption~\ref{a:prem}, and without it a belief inferred from no premises would form a tree with no roots; that is the shape of a priori justification, discussed in Section~\ref{s:limits}.

\subsection{The roots require no further support}

The roots are experiences, and the trilemma does not arise for them. The sceptic's demand is a demand for reasons, and reasons are what one can cite in support of a claim one makes. An experience makes no claim: it is not an assertion, a judgement or a commitment, and there is nothing that a subject who has an experience has done for which she can be asked to give reasons. She can be asked why she believes what the experience inclines her to believe, and she answers by citing the experience: \enquote{it looks red}. She cannot be asked why it looks red in the sense of a request for justification, though she can be asked in the sense of a request for a causal explanation, which is answered by the theory of perception rather than by epistemology. That is the sense in which the roots require no further support.

The point is not that experiences are self-justifying, which would assimilate them to Lewis's certain beliefs about the given; nor that they are infallible, which they are not, since an experience may present as true what is false. Experiences are fallible and yet not candidates for justification, as a thermometer reading is fallible and not a candidate for being justified. The fallibility of experience is accommodated in the tree by defeat: a belief grounded in an experience is justified only if undefeated, and a defeater may come from any other belief the subject holds.

\subsection{The isolation model}

The isolation objection to coherentism can be made precise. Coherentism counts as justified the members of a system each of which is supported by other members; formally, it takes the justified beliefs to be the largest class every member of which is supported by members of the class. Our account takes them to be the smallest class closed under the two sources. The two notions coincide when every belief in the system rests ultimately on experience, and diverge for isolated systems.

\begin{proposition}[Isolation]\label{p:isolation}
There is a system of two beliefs, each held by good inference from the other and neither grounded in experience, in which (i) every belief is supported by members of the system, so that the coherentist counts both as justified; and (ii) neither belief belongs to the smallest class closed under experiential grounding and good inference.
\end{proposition}

\begin{proof}
Let the beliefs be $a$ and $b$, with $a$ inferred from $b$ and $b$ from $a$, and nothing grounded. (i) is immediate. For (ii), suppose some belief belongs to the smallest closed class, and take a shortest derivation of one; it cannot end in grounding, since nothing is grounded, and if it ends in an inference, the premise has a shorter derivation, contradicting the choice.
\end{proof}

The model is the consistent novel in miniature. Its beliefs support each other and nothing else, and coherentism, in the absence of an observation requirement, must count them justified; Theorem~\ref{t:rootless} counts them unjustified. An observation requirement, such as BonJour's, makes the coherentist's notion coincide with ours, since it requires the system to answer to experience; that is the sense in which coherentism, so amended, is a tree theory under another name.

\section{Further cases}\label{s:cases}

\subsection{Memory}

A remembered belief, such as my belief that I parked on the third level, is not grounded in a present experience. It is justified, on the present view, because it preserves the justification it had when it was formed, by an experience that is no longer present. This is the preservationist account of memory \citep{Burge1993}. In the tree, memory is an edge that carries justification across time from a belief formed earlier to the same belief held now, and the root of the tree is the original experience. A memorial seeming, in which it seems to the subject that she remembers, may in addition confer prima facie justification by EPI-19; where the original experience never occurred, as in confabulation, the memorial seeming is the only support and is a root in its own right. Either way the tree is finite and ends in experience.

\subsection{Testimony}

EPI-08 holds that testimony transmits knowledge but does not generate it. The structure of justification explains why. When I believe what you tell me, my belief is held by inference from my belief that you told me so, together with background beliefs about your competence and sincerity, and my belief that you told me is grounded in my auditory experience of your utterance. But the tree does not end there. Testimony transmits the speaker's support, and its full support tree passes through the speaker's belief to the experiences on which that belief rests. The trees of testimonial justification thus span subjects: their roots lie, in part, in the experiences of others. \citet{Lackey2008} takes cases in which a hearer acquires knowledge that the speaker lacks to show that testimony can generate knowledge. On the diagnosis of EPI-08, what generates it in such cases is the hearer's own contribution, and the present structure agrees: the hearer's tree has roots of its own, in her experiences of the speaker's reliability and of the circumstances, which the speaker's tree lacks.

\subsection{Induction and rule-circularity}

EPI-03 argues that induction is justified inductively, the circularity being rule-circular rather than premise-circular. Theorem~\ref{t:rootless} excludes premise-circularity as a source of justification, and it is worth checking that it does not exclude the rule-circular argument. The argument for the reliability of induction takes as premises beliefs about the past successes of inductive inference, each of which is grounded in experience or derives from experience by good inference; its conclusion, that induction is reliable, is not among them. The tree is well founded and rooted in experience. What is circular is that the edge from premises to conclusion uses the inductive rule, whose goodness is what the conclusion asserts. Definition~\ref{d:infer} makes the goodness of a rule a property of the edge, not a node of the tree, and so the structure is indifferent to whether the rule is used to support a belief about itself. Whether the rule is in fact good is settled, as EPI-03 argues, by whether it is reliable in the world in which it is used, and \citet{Alston1986} shows that every basic source is in the same position. The tree structure therefore rules out exactly the circularity that EPI-03 concedes is vicious and permits exactly the circularity that EPI-03 argues is benign.

\subsection{Background beliefs and explanation}

Much empirical justification is by inference to the best explanation, which draws on background theory. A geologist who sees striations on a rock infers that a glacier passed over it. Her belief is held by good inference from her belief that there are striations, grounded in her experience, and from background beliefs about glacial action, each with its own tree. The background beliefs are many, but each is finitely supported, and the conclusion has a tree of height one greater than the highest of them. The appearance that explanatory inference is holistic, that the conclusion is supported by the whole of the geologist's theory at once, is the appearance created by the number of premises and by the role of the whole theory in defeat: a background belief that conflicted with glacial action would defeat the inference. Neither feature requires circular support.

\subsection{Expertise and perceptual learning}

An expert radiologist sees a tumour where a novice sees a shadow. The expert's belief is grounded in her experience, which, through training, presents the tumour, and she is justified by Assumption~\ref{a:exp}. The novice who believes that there is a tumour on the radiologist's say-so is justified by testimony. \citet{Siegel2017} argues that perceptual experiences can be formed rationally or irrationally, depending on the states that shape them, and that experiences shaped by prejudice or wishful thinking lose some of their justifying power. We think the cases are better described by defeat: the subject who learns that her experience has been shaped by a wish acquires an undercutting defeater, and EPI-19 treats them in that way. If Siegel is right, Assumption~\ref{a:exp} must be restricted to experiences formed without epistemic defect; the tree structure is unchanged, since the restriction concerns which experiences can be roots, not whether roots are experiences.

\section{Coherence in its place}\label{s:coherence}

Coherentists were right that coherence matters. The convergence of independent witnesses, the fit of an observation with a well-confirmed theory and the agreement of several instruments all increase justification. The question is how. \citet{CILewis1946} observed that the agreement of several independent and individually unreliable memories can make their common content highly credible, and argued that this effect requires that each memory have some initial credibility of its own: without it, congruence is worthless. The modern literature has made the point exact. \citet{KleinWarfield1994} showed that a more coherent set of beliefs can be less probable than a less coherent one. \citet{Shogenji1999} and \citet{Olsson2005} proved that coherence among reports increases the probability of their content only if the reports are independent and individually have some positive reliability, and \citet{BovensHartmann2003} proved that no measure of coherence is truth-conducive in general, other things being equal.

These results fit the tree structure. Coherence contributes to justification in two ways. First, it supplies premises: that several independent witnesses agree is itself a fact, believed on the basis of experiences of their testimony, and from it one may infer by good inference that what they agree on is true. The inference is a node with several children, each rooted in experience. Second, incoherence supplies defeaters: a belief that conflicts with others is defeated or defeats them. What coherence cannot do is generate justification for beliefs none of which has experiential support, which is what the isolation model shows and what the probabilistic results confirm. Haack's crossword is answered in the same way. The intersecting entries support one another because each is supported by its own clue, and the support an entry receives from its neighbours is a further inference from their clue-supported status, not a circle. A crossword with no clues, in which every entry is justified only by its fit with the others, has as many solutions as there are consistent fillings, and no entry in it is justified.

\section{Recent literature}\label{s:recent}

The last two decades have brought the regress debate into contact with probability theory and with the epistemology of perception. Three developments bear directly on the present account.

First, the probabilistic defence of infinitism. \citet{Peijnenburg2007} and \citet{AtkinsonPeijnenburg2017} showed that in an infinite chain in which each proposition is made probable by the next, the probability of the first can be well defined and determined by the conditional probabilities alone, the influence of distant links fading away. This answers the objection that infinite chains are formally incoherent. It does not answer Corollary~\ref{c:regress}. The probabilistic results concern the probability of the first member given the chain, and they show that, in the limit, the probability is independent of any starting point. But justification is not probability: an infinite chain of conditional probabilities that are not anchored to anything the subject has is a feature of the propositions, not of the subject's epistemic situation. The fading of distant links is, if anything, an argument for finite trees, since it shows that only the near links matter.

Second, the epistemology of perceptual seemings. \citet{Pryor2000} and \citet{Huemer2001} defended the view that perceptual experience gives immediate prima facie justification, and the view has become the main alternative to coherentism and reliabilism in the epistemology of perception. Modest foundationalists commonly describe it as a form of foundationalism with non-doxastic foundations \citep{Audi1993}. The present account agrees with the substance and disagrees with the label, and the difference is not merely verbal. A foundationalist with non-doxastic foundations still counts perceptual beliefs as basic, that is, as justified without the support of anything that itself has epistemic standing; and she inherits from classical foundationalism the question how basic beliefs can be justified without support, which leads to the demand for meta-justification pressed by \citet{BonJour1978}. On the tree account the perceptual belief has support, from the experience, and the question how a belief can be justified without support does not arise for any belief.

Third, the turn to McDowell's conceptualism. \citet{McDowell1994} argued that experience can be a reason for belief only if its content is conceptual, the same kind of content as belief; experience is passive, but it actualizes conceptual capacities. The view is designed to escape Sellars's dilemma, and it is compatible with the tree structure: if experiences have conceptual content, they can stand in rational relations to beliefs, and they remain roots, since their passivity means that they are not held for reasons. Our account does not need McDowell's conceptualism, but it can accept it.

\section{Objections}\label{s:objections}

\begin{objection}[Sellars's dilemma]
Either an experience has propositional content or it does not. If it does, it can be true or false, and the question whether one is entitled to it arises; it then needs justification like a belief. If it does not, it cannot stand in logical relations to beliefs and cannot justify them \citep{Sellars1956}.
\end{objection}

\begin{reply}
Experiences have propositional content, so the second horn does not apply. The first horn conflates two properties: having content that can be true or false, and being a state for which one can be asked for one's reasons. Beliefs have both. Experiences, like the readings of instruments, have the first and not the second: they can be accurate or inaccurate, but they are not adopted, and there is no entitlement to have them for which reasons could be given. A content that can be false can still justify, provided its possible falsity is accommodated by defeat, which is what Assumption~\ref{a:exp} requires.
\end{reply}

\begin{objection}[Only beliefs are reasons]
Nothing can count as a reason for holding a belief except another belief. Experiences cause beliefs, but a cause is not a reason \citep{Davidson1986}.
\end{objection}

\begin{reply}
The objection identifies reasons with premises. A belief grounded in an experience is not inferred from the experience, and the experience is not a premise; but it is the subject's reason in the sense that she holds the belief because of how things look and would cite how things look if asked. Davidson's alternative makes every reason a belief and so faces the isolation objection, which Proposition~\ref{p:isolation} makes exact. EPI-21 adds a further distinction: experiences are not evidence, since evidence consists of propositions known, but they confer justification. The relation of an experience to the belief it grounds is a distinct relation from the relation of evidence to hypothesis, and it is the one that makes contact with the world.
\end{reply}

\begin{objection}[Meta-justification]
If a perceptual belief is justified by an experience, the subject must have a reason to think that beliefs grounded in experiences of that kind are likely to be true. Otherwise the experience is, from her point of view, an arbitrary basis, and the belief is not justified for her \citep{BonJour1978}.
\end{objection}

\begin{reply}
The demand generates a regress of its own: the reason to think that experiences of that kind are reliable must itself be justified, ultimately by experiences whose reliability must in turn be vouched for. If the demand were correct, no belief would be justified. EPI-19 argues that the demand is unnecessary, since seemings confer prima facie justification without the subject's having a reason to take them to be reliable; and the companion argument of EPI-03 shows that the reliability of a source may be supported, once the source is used, by a rule-circular argument. The subject can have a reason to think her experiences reliable, and typically does; that reason is not a condition of their justifying.
\end{reply}

\begin{objection}[Arbitrariness]
Any stopping point in the chain of reasons is arbitrary: if the subject has no reason for her experience-based belief beyond the experience, she has no answer to the question why she should trust it, and her stopping there is dogmatic \citep{Klein1999}.
\end{objection}

\begin{reply}
The stopping point is not arbitrary, because it is not chosen. A subject who stops at a belief and declines to give reasons for it would be arbitrary; a subject whose chain of reasons ends in an experience has not stopped, because there is nothing further for her to give reasons for. The question \enquote{why trust your experience?} is a question about the reliability of a source, and it is answered, when it is answered, by an argument whose premises are themselves rooted in experience, as in the rule-circular case. Corollary~\ref{c:regress} shows, moreover, that the infinitist alternative generates no justification unless the chain makes contact with experience somewhere; and if it does, the infinite chain is idle. EPI-05-b develops this objection.
\end{reply}

\begin{objection}[Mutual support]
Some support is irreducibly mutual. In a crossword, each entry is supported by the entries that cross it, and they by it; in science, theory and observation support each other. A tree cannot represent this \citep{Haack1993}.
\end{objection}

\begin{reply}
A tree can represent the effect of mutual support, as Section~\ref{s:coherence} argued. Where $a$ and $b$ each make the other more credible, each has its own experiential tree, and the belief in $a$ may in addition be supported by an inference from the belief in $b$ held on its own tree, and vice versa. The resulting structure, unfolded, is a tree of finite height in which each support relation is directed from a belief supported independently to one it further supports. What a tree cannot represent is support that exists only in the circle, and Theorem~\ref{t:rootless} shows that there is no such support.
\end{reply}

\section{What the result does not establish}\label{s:limits}

The tree theorem establishes the structure of justification given the three assumptions about its sources. It does not establish which inferences are good or which experiences can be roots; those questions belong to EPI-03, EPI-19 and the epistemology of particular domains. It concerns empirical justification. EPI-12 holds that there is a priori knowledge of logic, arithmetic and some conceptual truths. In the present framework, logical and arithmetical principles figure in empirical trees mainly as rules of inference, that is, as edges rather than nodes. Where they are believed and justified, their justification has a structure that the present manuscript does not determine: an inference from no premises, which Assumption~\ref{a:prem} excludes for empirical beliefs, or a root in an intellectual seeming, which EPI-19's principle covers but which is not a perceptual experience. Finally, the theorem concerns propositional justification and its structure. It does not address doxastic justification, which requires in addition that the belief be based on its support, or knowledge, which on EPI-01 requires that the belief be safe.

\section{Formalization}

The module \nolinkurl{Philosophy/EPI/EPI_05.lean} represents beliefs and experiences as distinct primitive types, grounding, good inference, defeat and justification as primitive relations and predicates, and Assumptions~\ref{a:exp}--\ref{a:prem} as axioms. Support trees are an inductive predicate following Definition~\ref{d:tree}. It proves Lemmas~\ref{l:mono} and~\ref{l:bound}, Theorems~\ref{t:tree}, \ref{t:rootless} and~\ref{t:dep}, Corollaries~\ref{c:circle} and~\ref{c:regress}, and Proposition~\ref{p:isolation}, the last in a concrete model with no axioms. The main statements are:

\begin{lstlisting}
theorem tree_theorem (b : Belief) :
    Justified b ↔ ∃ n, Tree n b

theorem no_rootless (S : Belief → Prop)
    (hg : ∀ e b, S b → ¬ Grounds e b)
    (hi : ∀ ps b, S b → Infers ps b → ∃ p, p ∈ ps ∧ S p) :
    ∀ b, S b → ¬ Justified b

theorem experiential_dependence (b : Belief) (hj : Justified b) :
    ∃ e, DerivesFrom b e
\end{lstlisting}

The \texttt{\#print axioms} output, with the namespace prefix omitted, is:

\begin{lstlisting}
'tree_theorem' depends on axioms:
  [Belief, Defeated, Experience, Grounds, Infers, Justified,
   experiential, no_other_source, transmission]
'no_rootless' depends on axioms:
  [Belief, Defeated, Experience, Grounds, Infers, Justified,
   no_other_source]
'experiential_dependence' depends on axioms:
  [Belief, Defeated, Experience, Grounds, Infers, Justified,
   no_other_source, premises_nonempty]
'Isolation.not_generated' does not depend on any axioms
\end{lstlisting}

\begin{sloppypar}
\texttt{Belief}, \texttt{Experience}, \texttt{Grounds}, \texttt{Infers}, \texttt{Defeated} and \texttt{Justified} are the primitives of Definitions~\ref{d:ground}--\ref{d:defeat}; \texttt{experiential} is Assumption~\ref{a:exp}, \texttt{transmission} Assumption~\ref{a:trans}, \texttt{no\_other\_source} Assumption~\ref{a:nos} and \texttt{premises\_nonempty} Assumption~\ref{a:prem}. The formalization brings out a division of labour among the assumptions that the prose proofs leave implicit. Theorem~\ref{t:rootless}, and with it the rejection of circles and regresses, uses only Assumption~\ref{a:nos}: the negative half of the view follows from the claim that justification has no sources other than experience and inference, whatever those sources deliver. Theorem~\ref{t:tree} uses in addition the two sufficiency assumptions, which are needed for the direction from trees to justification. Theorem~\ref{t:dep} uses Assumption~\ref{a:nos} and Assumption~\ref{a:prem} and not the sufficiency assumptions. The finiteness of trees is not assumed anywhere; it is derived from the finiteness of premise lists, through Lemma~\ref{l:bound}. No classical axiom is used.
\end{sloppypar}

\section{Conclusion}

Justification is neither a pyramid resting on basic beliefs nor a raft of mutually supporting ones, nor an infinite chain. It is a tree. Every justified belief is supported, and its support, traced downwards, passes through finitely many good inferences to beliefs grounded in perceptual experiences, which are not beliefs and are not candidates for justification. Given that justification is generated by undefeated experiential grounding and transmitted by good inference, and has no other sources, a belief is justified exactly when it has such a tree; circles and infinite regresses generate nothing; every justified empirical belief derives from experience; and coherence contributes by supplying premises and defeaters, not by generating justification on its own.

\begin{thebibliography}{99}\small
\bibitem[Alston(1986)]{Alston1986} Alston, W.~P. (1986). Epistemic circularity. \emph{Philosophy and Phenomenological Research}, 47(1), 1--30.
\bibitem[Aristotle(c.~350 BCE)]{AristotlePostAn} Aristotle (c.~350 BCE). \emph{Posterior Analytics}.
\bibitem[Atkinson and Peijnenburg(2017)]{AtkinsonPeijnenburg2017} Atkinson, D. and Peijnenburg, J. (2017). \emph{Fading Foundations: Probability and the Regress Problem}. Springer.
\bibitem[Audi(1993)]{Audi1993} Audi, R. (1993). \emph{The Structure of Justification}. Cambridge University Press.
\bibitem[BonJour(1978)]{BonJour1978} BonJour, L. (1978). Can empirical knowledge have a foundation? \emph{American Philosophical Quarterly}, 15(1), 1--14.
\bibitem[BonJour(1985)]{BonJour1985} BonJour, L. (1985). \emph{The Structure of Empirical Knowledge}. Harvard University Press.
\bibitem[BonJour and Sosa(2003)]{BonJourSosa2003} BonJour, L. and Sosa, E. (2003). \emph{Epistemic Justification: Internalism vs.\ Externalism, Foundations vs.\ Virtues}. Blackwell.
\bibitem[Bovens and Hartmann(2003)]{BovensHartmann2003} Bovens, L. and Hartmann, S. (2003). \emph{Bayesian Epistemology}. Oxford University Press.
\bibitem[Burge(1993)]{Burge1993} Burge, T. (1993). Content preservation. \emph{The Philosophical Review}, 102(4), 457--488.
\bibitem[Davidson(1986)]{Davidson1986} Davidson, D. (1986). A coherence theory of truth and knowledge. In E.~LePore (Ed.), \emph{Truth and Interpretation: Perspectives on the Philosophy of Donald Davidson} (pp.~307--319). Blackwell.
\bibitem[Descartes(1641)]{Descartes1641} Descartes, R. (1641). \emph{Meditationes de prima philosophia}.
\bibitem[Haack(1993)]{Haack1993} Haack, S. (1993). \emph{Evidence and Inquiry: Towards Reconstruction in Epistemology}. Blackwell.
\bibitem[Huemer(2001)]{Huemer2001} Huemer, M. (2001). \emph{Skepticism and the Veil of Perception}. Rowman \& Littlefield.
\bibitem[Klein(1999)]{Klein1999} Klein, P.~D. (1999). Human knowledge and the infinite regress of reasons. \emph{Philosophical Perspectives}, 13, 297--325.
\bibitem[Klein(2007)]{Klein2007} Klein, P.~D. (2007). Human knowledge and the infinite progress of reasoning. \emph{Philosophical Studies}, 134(1), 1--17.
\bibitem[Klein and Warfield(1994)]{KleinWarfield1994} Klein, P. and Warfield, T.~A. (1994). What price coherence? \emph{Analysis}, 54(3), 129--132.
\bibitem[Lackey(2008)]{Lackey2008} Lackey, J. (2008). \emph{Learning from Words: Testimony as a Source of Knowledge}. Oxford University Press.
\bibitem[Lewis(1946)]{CILewis1946} Lewis, C.~I. (1946). \emph{An Analysis of Knowledge and Valuation}. Open Court.
\bibitem[McDowell(1994)]{McDowell1994} McDowell, J. (1994). \emph{Mind and World}. Harvard University Press.
\bibitem[Neurath(1932)]{Neurath1932} Neurath, O. (1932). Protokolls\"atze. \emph{Erkenntnis}, 3, 204--214.
\bibitem[Olsson(2005)]{Olsson2005} Olsson, E.~J. (2005). \emph{Against Coherence: Truth, Probability, and Justification}. Oxford University Press.
\bibitem[Peijnenburg(2007)]{Peijnenburg2007} Peijnenburg, J. (2007). Infinitism regained. \emph{Mind}, 116(463), 597--602.
\bibitem[Pryor(2000)]{Pryor2000} Pryor, J. (2000). The skeptic and the dogmatist. \emph{No\^us}, 34(4), 517--549.
\bibitem[Schlick(1934)]{Schlick1934} Schlick, M. (1934). \"Uber das Fundament der Erkenntnis. \emph{Erkenntnis}, 4, 79--99.
\bibitem[Sellars(1956)]{Sellars1956} Sellars, W. (1956). Empiricism and the philosophy of mind. \emph{Minnesota Studies in the Philosophy of Science}, 1, 253--329.
\bibitem[Sextus Empiricus(c.~200)]{Sextus} Sextus Empiricus (c.~200). \emph{Outlines of Pyrrhonism}.
\bibitem[Shogenji(1999)]{Shogenji1999} Shogenji, T. (1999). Is coherence truth conducive? \emph{Analysis}, 59(4), 338--345.
\bibitem[Siegel(2017)]{Siegel2017} Siegel, S. (2017). \emph{The Rationality of Perception}. Oxford University Press.
\bibitem[Sosa(1980)]{Sosa1980} Sosa, E. (1980). The raft and the pyramid: Coherence versus foundations in the theory of knowledge. \emph{Midwest Studies in Philosophy}, 5, 3--25.
\end{thebibliography}

\end{document}
